%! Author = lili
%! Date = 4/27/26

\untit{Vorlesung 25.04.2025 (V02 2025), ergänzt mit dem Skript aus 2026}

\section*{Lernziele}
\begin{itemize}
    \item Unterschied zwischen disjunkten und paarweise disjunkten Mengen
    \item Rechenregeln für Wahrscheinlichkeiten
    \begin{itemize}
        \item Komplementär Wahrscheinlichkeit
        \item Ein- und Ausschlussprinzip
        \item Spezialfall paarweise disjunkter Mengen
    \end{itemize}
\end{itemize}

%%%%%%%%%%%%%%%%%%%%%
% Beispiele Wraum %
%%%%%%%%%%%%%%%%%%%%%

\section{Wahrscheinlichkeitsräume, deren Ereignisraum nicht abzählbar ist}\label{sec:wahrscheinlichkeitsraume-deren-ereignisraum-nicht-abzahlbar-ist}
\begin{bsp}
    \textbf{\gls{wraum}}  \txc{(2025)} \\
    Stoppen einer Zeit bis zu einer Stunde in ms.
    \begin{align*}
        \eraum  & = \{0,1,\dots,360 000\} & \text{(Zeit in ms)} \\
        \glssymbol{wfk}(\glssymbol{elementarereignis}) & \coloneqq \frac{1}{\# \eraum } \quad \forall \glssymbol{elementarereignis} \in \eraum  & \text{(Gleichverteilung)} \\
        \pe(\erga) & = \frac{\# \erga}{\# \eraum } \\
        \erga & = \{600 000, \dots , 9000 000\} & \text{Zeit liegt zw. 10 und 15 Minuten} \\
        \pe(\erga) & = \frac{300 001}{3 600 001} \approx \frac{300 000}{3 600 000} = \frac{1}{12}
    \end{align*}
    Wir wollen kontinuierlich modellieren:
    \begin{align*}
        \tilde {\eraum}  & = [0,60] & \text{(in Minuten)}\\
        \text{bzw. }  \tilde {\eraum}  & = [0, 360 000]  & \text{(in ms)} \\
        \tilde \erga & = [10,15] \\
        \mathbb{\tilde P} (\tilde \erga) & \overset{!}{=} \frac{15-10}{60-0} = \frac{5}{60} = \frac{1}{12}
    \end{align*}
    Um auch modellieren zu können, dass Intervalle \gls{ws} haben, die nicht \texttt{iwas} von ihrer Länge
    abgängen, arbeiten wir mit R-Integralen. \texttt{R-Integral?}
\end{bsp}
\komm{TODO: Was genau ist $\tilde {\eraum} $ ? } % TODO

%%%%%%%%%%%%%%%%%%%%%
% Dichten %
%%%%%%%%%%%%%%%%%%%%%
\subsection{Eindimensionale Dichten}\label{subsec:eindimensionale-dichten}
\begin{defi}
    \label{def:dichte}
    \textbf{Eindimensionale \dichte} \txc{(2025, Skript 26 - Definition 1.22)} \\
    EIne Funktion $\dcht : \mathbb{R} \rightarrow [0, \infty)$, die (uneigentlich) Riemann-integrierbar ist, heißt
    \underline{\dichte}, sofern
    \[
        \int_{- \infty}^{\infty}\dcht  (x) dx = 1
    \]
\end{defi}

\begin{bem}
    \defirefman{dichte}  \txc{(Skript 26 - Bemerkung 1.23)}

    \paragraph{Dichten sind nichtnegativ:} Die Definition einer \dichte f beinhaltet die Forderung
    $\dcht \leq 0$, das heißt,
    es muss
    \[
        \dcht (x) \geq 0 \forall x \in \rz
    \]
    gelten.
\end{bem}
Die folgende Bemerkung erinnert uns wohl an die für den Zweck dieser Vorlesung wichtigsten Aussagen über
die Existenz von Riemann-Integralen.
\begin{bem}
    \defirefman{dichte}  \txc{(2025, Skript 26 - Bemerkung 1.24)}

    \paragraph{Riemann-Integrierbarkeit}
    \begin{itemize}
        \item  % a
        \item $\dcht$ monoton oder $\dcht$ stetig $\Rightarrow \int_a^b \dcht(x) dx$ existiert für alle $(a, b \in \mathbb{R})$
        \txc{auch: für $- \infty < a < b < \infty$}
        % b
        \item $\dcht$ \gls{stst} oder stückweise monoton, $\Rightarrow \int_a^b \dcht(x) dx$ existiert
        $(a, b \in \mathbb{R})$
        \txc{auch: für $- \infty < a < b < \infty$}
        % e
        \item $\dcht \geq 0:\dcht $ Riemann-integrierbar
        $\Leftrightarrow \underset{N \rightarrow \infty}{\lim} \int_{-N}^N \dcht(x) dx < \infty$
        \begin{itemize}
            \item Damit ist gemeint: Ist $\dcht \geq 0$, so ist $\dcht$ genau dann uneigentlich Riemann-integrierbar über
            $\rz$, wenn das Riemann-Integral $\int_a^b \dcht(x) dx$ für jede Wahl $- \infty < a < b < \infty$
            existiert und $\underset{N \rightarrow \infty}{\lim} \int_{-N}^N \dcht(x) dx$
            existiert.
            In diesem Fall gilt $\int_{- \infty}^{\infty} \dcht(x) dx = \underset{N \rightarrow \infty}{\lim} \int_{-N}^N \dcht(x) dx$
        \end{itemize}
        % d
        \item $\dcht$ ist uneigentlich Riemann-integrierbar über R, gdw.
        sowohl $(- \infty, 0] $ als auch $[0, \infty)$
        uneigentlich Riemann-integrierbar ist.
        % c TODO
        \item $\dcht$ ist genau dann uneigentlich Riemann-integrierbar über $[0,\infty)$, wenn $\dcht$ über
        jedes Intervall $[0,b]$ Riemann-integrierbar ist und $\lim_{b \rightarrow \infty} \int_0^b \dcht(x) \dx$
        existiert.
        In diesem Fall schreiben wir
        \[
            \int^{\infty}_0 \dcht (x) \dx =  \lim_{b \rightarrow \infty} \int_0^b \dcht(x) \dx
        \]
        \item Analog: $\dcht$ ist genau dann uneigentlich Riemann-integrierbar über $(-\infty, 0]$, wenn $\dcht$
        über jedes Intervall $[a,0]$ Riemann-integrierbar ist $\lim_{a \rightarrow -\infty} \int_a^0 \dcht(x) \dx$ existiert.
        In diesem Fall schreiben wir
        \[
            \int_{-\infty}^0 \dcht (x) \dx =  \lim_{a \rightarrow - \infty} \int_a^0 \dcht(x) \dx
        \]
    \end{itemize}
\end{bem}

\karkar{def_dichte}
\karkar{bem_dichte_nichtnegativ}
\karkar{bem_riemann_integrierbarkeit_kriterien}
\karkar{bem_uneigentliche_integrierbarkeit_kriterium}
\karkar{bem_integrierbarkeit_halbachsen}

%%%%%%%%%%%%%%%%%%%%%
% Wahrscheinlichkeitsräume mit Dichten %
%%%%%%%%%%%%%%%%%%%%%
\subsection{Wahrscheinlichkeitsräume mit Dichten}
\begin{defi}
    \textbf{Wahrscheinlichkeitsräume mit Dichten}\label{def:wr_dichte} \txc{(2025, Skript 26 - Definition 1.25)} \\
    Sei eine \dichte $\dcht \geq 0$ gegeben (also $f: \rz \rightarrow [0,\infty)$.
    Wir setzen
    \begin{align*}
        \pe([a,b]) & \coloneqq  \int^b_a \dcht(x) dx & \text{für alle } - \infty < a < b < \infty \\
        \pe([-\infty, b]) & \coloneqq  \int_{-\infty}^b \dcht(x) dx  \\
        \pe([a, \infty]) & \coloneqq  \int_a^{\infty} \dcht(x) dx \\
        \pe(\bigcup_{i \in I} [a_i, b_i]) & \coloneqq  \sum_{i \in I} \int^{b_i}_{a_i} \dcht(x) dx &
        \text{für $I \subseteq\mathbb{N}$ (abzählbar)  und $a_i < b_i \leq a_{i+1}\txc{*}$}
    \end{align*}
    Falls
    \[
        a_1 = - \infty : (- \infty, b_1]\quad ; \ergb =+ \infty : [a, \infty)
    \]
    \txc{* Auch geschrieben als: Seien
        $I \subseteq\mathbb{N}$ eine abzählbare Indexmenge und $[a_i, b_i]$ Intervalle mit
        $a_i \leq b_i \leq a_{i+1} \leq b_{i+1}$ für
        alle $i, i + 1 \in I$. }
\end{defi}
\begin{bsp}
    \txc{(2025)}
    \texttt{fehlt}
\end{bsp}

\komm{\defirefman{wr_dichte} Die Definition von $\pe$ auf einer abzählbaren Vereinigung disjunkter,
    aufsteigend geordneter Intervalle $[a_i,b_i]$ ist additiv über die Dichte
    $\dcht$ definiert. Uneigentliche Randintervalle ($a_1=-\infty$ bzw.
    $b_{\max}=+\infty$) werden gesondert als Halbintervalle behandelt.
    Eine alternative, schwächere Formulierung erlaubt statt $a_i < b_i$ auch
    $a_i\leq b_i$, wodurch entartete (einpunktige) Intervalle zugelassen
    werden; da diese jeweils das Maß $0$ besitzen, ändert sich der Wert von
    $\pe$ dadurch nicht.}

\begin{bem}
    \defirefman{wr_dichte} \txc{(Skript 26 - Bemerkung 1.26)}
    \begin{itemize}
        \item Wann immer das Wahrscheinlichkeitsmass $\pe$ durch eine eindimensionale Dichte gegeben ist,
        betrachten wir nur Wahrscheinlichkeiten von Mengen, die wir darstellen können als maximal
        abzählbar unendliche Vereinigung oder als maximal abzählbar unendlichen Schnitt von Intervallen.
        Auch können wir wiederum Vereinigungen und Schnitte von derart erhaltenen Mengen
        betrachten.
        Diese Prozedur kann beliebig oft wiederholt werden.
        Trotzdem gibt es Mengen, denen wir keine Wahrscheinlichkeit zuordnen können.
        Da es nicht
        ganz einfach ist, derartige - pathologische - Mengen zu konstruieren, bedeutet das für uns keine
        wesentliche Einschränkung.
        \item Um ein Zufallsexperiment zu modellieren, in dem nur Elementarereignisse aus einer echten
        Teilmenge $I \in \rz$ angenommen werden können, wählen wir eine Dichte, die außerhalb von I
        Null ist, das heißt, es soll $\dcht (x) = 0$ für alle $x \notin I$
    \end{itemize}
\end{bem}
% Im folgenden Beispiel 1.30 werden wir
%diese Idee zum ersten Mal anwenden.

\begin{lem}
    \textbf{Wahrscheinlichkeit einelementiger Mengen} \txc{(Skript 26 - Lemma 1.27)} \\
    Ist ein Wahrscheinlichkeitsmass $\pe$ durch eine \dichte f gegeben, so gilt
    \[
        \pe(\{a\}) = 0 \quad \forall a \in \rz
    \]
    \paragraph{Beweis.} Skript 26 Seite 11 % TODO
\end{lem}

\begin{kor}
    \textbf{Dichten sind nicht eindeutig bestimmt} \txc{(Skript 26 - Korollar 1.28)} \\
    Das Wahrscheinlichkeitsmass $\pe$ sei durch eine Dichte f gegeben.
    Sei g eine weitere Dichte, die sich in
    höchstens abzählbar vielen Punkten von f unterscheidet, das heißt, es gebe eine abzählbare Menge M
    derart, dass $g(x) = f (x)$ für alle $x \in \rz \backslash M$.
    Dann definiert g dasselbe Wahrscheinlichkeitsmass $\pe$
    wie f.
\end{kor}

\begin{kor}
    \textbf{Wahrscheinlichkeit von Intervallen} \txc{(Skript 26 - Korollar 1.29)} \\
    Sei $\pe$ durch eine Dichte f gegeben.
    Dann gilt für jede Wahl von $a,b \in \rz$ mit $- \infty < a < b < \infty$,
    \[
        \pe([a, b]) = \pe((a, b]) = \pe([a, b)) = \pe((a, b))
    \]
\end{kor}

\begin{bsp}\label{bsp:1-30}
    \textbf{Messen einer zufälligen Zeit} \txc{(Skript 26 - Beispiel 1.30)} \\
    Wir wählen $\eraum = \rz$ und die Dichte
    \[
        \dcht(x) = \begin{cases}
                       \frac{1}{60} & x \in [0,60] \\
                       0 & sonst
        \end{cases}
    \]
    Dabei verstehen wir $\oio$ als die gemessene Zeit (in Minuten) bei einer Realisierung des Zufallsexperiments.
    Die Tatsache, daß $f = 0$ außerhalb von $[0, 60]$ gilt, reflektiert, daß in unserem Experiment
    nur Zeiten innerhalb einer Stunde ab Start möglich sein sollen. \\
    Wir müssen zeigen, daß f eine Dichte ist
    \begin{enumerate}
        \item $f \geq 0$ gilt per Definition von f
        \item f ist \gls{stst} und daher Riemann-integrierbar
        \item Die Normierungsbedingung ist erfüllt, weil
        \[
            \intinfs \dcht(x) \dx = \int_0^{60} \dcht(x) \dx = \int_0^{60} \frac{1}{60} \dx = \frac{60 - 0}{60} = 1
        \]
    \end{enumerate}
    Wir sehen, dass, wie zuvor,
    \[
        \pe([10, 15]) = \int_{10}^{15}  \dcht(x) \dx \int_{10}^{15} \frac{1}{60} \dx = \frac{5}{60} = \frac{1}{12}
    \]
    gilt.
\end{bsp}

\begin{bem}
    \textbf{Vorgehen beim Berechnen von Integralen mit stückweise gegebener Dichte}
    \txc{(Skript 26 - Bemerkung 1.31)} \\
    Es empfiehlt sich, sukzessive wie in obigem Beispiel vorzugehen:
    \begin{enumerate}
        \item Integral über Intervall aufschreiben
        \item Grenzen anpassen, so daß nicht mehr über Intervalle integriert wird, auf denen $f = 0$ gilt
        \item Konkrete Form von f einsetzen
        \item Integral ausrechnen
    \end{enumerate}
    Genauer
    \begin{enumerate}
        \item Integral über Intervall aufschreiben:
        \begin{itemize}
            \item $\intinfs \dcht(x) \dx$, um zu verifizieren, daß f die Normierungsbedingung erfüllt
            \item $\pe([a, b]) = \intinfs \dcht(x) \dx$, um eine Wahrscheinlichkeit zu berechnen
        \end{itemize}
        \item Grenzen anpassen, so daß nicht mehr über Intervalle integriert wird, auf denen $f = 0$ gilt:
        \begin{itemize}
            \item $\intinfs \dcht(x) \dx = \int_0^{60} \dcht(x) \dx$
            \item $\pe([50, 70]) = \int_{50}^{70} \dcht(x) \dx  = \int_{50}^{60} \dcht(x) \dx $
        \end{itemize}
        \item Konkrete Form von f einsetzen:
        \begin{itemize}
            \item $ \pe([10, 15]) = \int_{10}^{15}  \dcht(x) \dx \int_{10}^{15} \frac{1}{60} \dx$
            \item $\pe([50, 70]) = \int_{50}^{70} \dcht(x) \dx  = \int_{50}^{60} \dcht(x) \dx =  \int_{50}^{60} \frac{1}{60} \dx$
        \end{itemize}
        \item Integral ausrechnen:
        \begin{itemize}
            \item $ \pe([10, 15]) = \int_{10}^{15}  \dcht(x) \dx \int_{10}^{15} \frac{1}{60} \dx = 1$
            \item $\pe([50, 70]) = \int_{50}^{70} \dcht(x) \dx  = \int_{50}^{60} \dcht(x) \dx =  \int_{50}^{60} \frac{1}{60} \dx = \frac{1}{6}$
        \end{itemize}
    \end{enumerate}
\end{bem}

\begin{bem} \txc{(Skript 26)}
Falls f durch eine Fallunterscheidung mit mehr Fällen als in unserem Beispiel gegeben ist, so muss im
    dritten Schritt das Integral entsprechend aufgespalten werden,
    \[
        \int^d_a \dcht(x) \dx = \int^b_a \dcht(x) \dx + \int^c_b \dcht(x) \dx + \int^d_c \dcht(x) \dx
    \]
    Für $a < b < c < d$, bevor die konkrete Form von f einsetzt wird. \\
    Losgelöst von der Idee, daß wir eine zufällige Zeit messen wollen, können wir ganz allgemein für ein
    beliebiges beschränktes Intervall I eine Dichte definieren derart, daß die Wahrscheinlichkeit, daß der
    Ausgang des entsprechenden Zufallsexperiments in einem Teilintervall $J \subseteq I$ liegt, weiterhin nur von
    der Länge des entsprechenden Teilintervalls J abhängt
\end{bem}

\karkar{def_wraum_mit_dichte}
\karkar{bem_grenzen_der_wsmodellierung_mit_dichte}
\karkar{bem_dichte_null_ausserhalb_I}
\karkar{lem_einelementige_mengen_wahrscheinlichkeit_null}
\karkar{kor_dichten_nicht_eindeutig}
\karkar{kor_ws_intervalle_rand_egal}
\karkar{bem_vorgehen_integrale_stueckweise_dichte}
\karkar{bem_integral_aufspalten_mehr_faelle}

%%%%%%%%%%%%%%%%%%%%%
% Uniforme Verteilung %
%%%%%%%%%%%%%%%%%%%%%
\subsection{Uniforme Verteilung}
\begin{defi}
    \textbf{Uniforme Verteilung auf einem Intervall}
    \label{def:uniform_intervall} (Kontinuumsäquivalent zu gleichverteilt)
    \txc{(2025, Skript 26 - Definition 1.32)} \\
    Die uniforme Verteilung auf einem beschränkten Intervall $I \subseteq \mathbb{R}$ (äquivalent zu $I = [a, b]$ mit $- \infty < a < b < \infty$) ist gegeben
    durch eine \dichte
    \[
        \dcht(x) = \begin{cases}
                       \frac{1}{\mid I \mid } & x \in I \\ 0 & sonst
        \end{cases} \quad = \begin{cases}
                                                   \frac{1}{b-a} & x \in I \\ 0 & sonst
        \end{cases}
    \]
    wobei $\mid I \mid $ die Intervalllänge bezeichnet. \\
\end{defi}

\begin{bsp} \defirefman{uniform_intervall} txc{(Skript 26 - Definition 1.32)} \\
    Sei nun $[a, b] = [0, 10]$.
    Die Wahrscheinlichkeit, daß der Ausgang des Zufallsexperiments sehr klein oder sehr groß ist, sagen wir entweder kleiner als 0.1 oder größer als 9.9 ist, beträgt
    \[
        \pe((- \infty, 0.0) \cup (9.9, \infty)) = \int_0^{0.1} \dcht(x) \dx + \int_{9.9}^{10} \dcht(x) \dx = \frac{1}{10} [0.1 + 0.1] = \frac{2}{100} = \frac{1}{50}
    \]
\end{bsp}

\begin{bem} \bspref{1-30}
    \txc{(Skript 26 - Bemerkung 1.33)} \\
Der Wert $\frac{1}{b-a}$ ergibt sich aus der
Normierungsbedingung.
Daraus ergibt sich, daß es auf einem nicht beschränkten Intervall keine uni-
forme Verteilung geben kann.
Der Beweis verwendet dieselbe Idee wie jene, die wir in Bemerkung \autoref{bem:1-18}
benutzt haben.
\end{bem}
Wir können die Definition der uniformen Verteilung verallgemeinern.
\begin{defi}
    \textbf{Uniforme Verteilung auf einer Menge M} \txc{(Skript 26 - Definition 1.34)} \\
    Sei $M \subseteq \rz$ eine Menge, der wir einen Inhalt $|M| = \lambda(M )$ zuordnen können.
    Ist
    $|M| < \infty$, so ist die
    uniforme Verteilung auf M gegeben durch die Dichte
    \[
        \dcht(x) = \begin{cases}
                       \frac{1}{|M|} & x \in M \\
                       0 & sonst
        \end{cases}
    \]
\end{defi}

\begin{bem}
    \txc{(Skript 26 - Bemerkung 1.35)} \\
    Die uniforme Verteilung auf einer beschränkten Menge $M \subseteq \rz$ ist das kontinuierliche Analogon der
    Gleichverteilung auf einer endlichen Menge.
    Es gilt
    \[
        \pe(A)  =\frac{|A \cap M|}{|M|}
    \]
    für alle $A \subseteq \rz$, für die der Inhalt $|A|$ definiert ist. \\
    Wartezeiten werden häufig mit Hilfe der Exponentialverteilung modelliert. % TODO Karteikarte
\end{bem}

\karkar{def_uniforme_verteilung_intervall}
\karkar{bem_uniform_nur_beschraenkt}
\karkar{def_uniforme_verteilung_menge_M}
\karkar{bem_uniform_M_kontinuierliches_analogon}

%%%%%%%%%%%%%%%%%%%%%
% Exponentialverteilung %
%%%%%%%%%%%%%%%%%%%%%
\subsection{Exponentialverteilung}
\begin{defi}
    \textbf{Exponentialverteilung} \label{def:exponentialverteilung} \txc{(2025, Skript 26 - Definition 1.36)}\\
    Die Exponentialverteilung zum Parameter $\alpha > 0$ ist gegeben durch die \dichte
    \begin{equation}
        \label{eq:Exponentialverteilung}
        \dcht(x) = \begin{cases}
                       0 & \text{ für } x < 0 \\
                       \alpha e ^{- \alpha x} &  \text{ für } x \geq 0
        \end{cases}
    \end{equation}
\end{defi}
\begin{bem}
    \defirefman{exponentialverteilung} \textbf{ - Die Exponentialverteilung ist wohldefiniert}  \txc{(2025, Skript 26 - Lemma 1.37)} \\
    Gleichung~\ref{eq:Exponentialverteilung} ist eine \dichte:
    \begin{itemize}
        \item $\dcht  \geq 0 $
        \item $\dcht$ R-integrierbar, da \gls{stst}
        \orangenote{$g(x) = e^{-\alpha x}$ \\ $g'(x) = (-\alpha) e^{- \alpha x}$ (Produktregel)}
        \item \texttt{???}:
        \begin{align*}
            \int_{-\infty}^{\infty} \dcht(x) dx & =   \int_{0}^{\infty} \dcht(x) dx  \\
            & =   \int_{0}^{\infty} \alpha e^{-\alpha x} dx  \\
            & = \underset{N \rightarrow \infty}{lim} \int^N_0 \alpha e ^{-\alpha x} dx \\
            & = \underset{N \rightarrow \infty}{lim} [-e ^{-\alpha x} \mid^N_{x=0}] \\
            & = \underset{N \rightarrow \infty}{lim} [-e ^{-\alpha N}
            \txo{\underbrace{
                \textcolor{black}{- ( -\txo{\underbrace{
                    \textcolor{black}{e^{-\alpha \cdot 0}
                    }}_{=1}}}
            }_{=1}}
            )] \\
            & = 1 - \txo{
                \underbrace{
                    \textcolor{black}{\underset{N \rightarrow \infty}{lim} e ^{- \alpha N}}
                }_{=1}
            } \\
            & = 1
        \end{align*}
    \end{itemize}
\end{bem}



\chapter{Rechnen mit Wahrscheinlichkeiten (V02)}
\label{ch:rechnen-mit-wahrscheinlichkeiten}
\untit{Vorlesung 25.04.2025 (V02 2025), ergänzt durch Skript 2026 für Vorlesung 24.04.2026 (V02 2026)}

\section*{Lernziele}
\begin{itemize}
    \item Unterschied zwischen disjunkten und paarweise disjunkten Mengen kennen
    \item Rechenregeln für Wahrscheinlichkeiten anwenden und beweisen können, insbesondere
    \begin{itemize}
        \item Komplementärwahrscheinlichkeit (``Gegenwahrscheinlichkeit'')
        \item Ein-und-Ausschlußprinzip
        \item Spezialfall: Wahrscheinlichkeit der Vereinigung paarweise disjunkter Mengen
    \end{itemize}
\end{itemize}

\section{Mengenoperationen und Ereignisse}\label{sec:mengenoperationen-und-ereignisse}
\begin{defi}\label{def:komplement} \textbf{Komplement von Mengen} \txc{(Skript 2026 - Definition 2.1)}) \\
    Gilt $\erga \subseteq \eraum$, so schreiben wir $\erga^c$ für das Komplement von $\erga$ in $\eraum$, das heißt, wir setzen
    \[
        \erga^c = \eraum \backslash \erga
    \]
\end{defi}
\begin{bem}
    \defirefman{komplement} \txc{(Skript 2026 - Definition 2.1)}) \\
    Diese Notation ist mit Vorsicht anzuwenden, da sie voraussetzt, dass klar ist, auf welche Referenz-
    Obermenge $\eraum$ wir uns beziehen.
\end{bem}
\begin{defi}
    \textbf{Vereinigung/Schnitt} abzählbar vieler Mengen \txc{(2025, Skript 2026 - Definition 2.2)} \\
    $I \subseteq \mathbb{N}$ Indexmenge. \\
    Für jedes $i \in I $ sei $\erga_i \subseteq \eraum$:
    \begin{itemize}
        \item $\erga_i \in \mathbb{R}(\eraum )$ für $\eraum  $ abzählbar
        \item $\erga_i$ ein Intervall für $\eraum  = \mathbb{R}$ (oder $\eraum  = [0, \infty)$ oder $\eraum  = [a,b]$)
    \end{itemize}
    \begin{itemize}
        \item[(a)] $\underset{i \in I}{\bigcup} \erga_i = \{\glssymbol{elementarereignis} \in \eraum  : \glssymbol{elementarereignis} \in \erga_i \text{ für mind. ein } i \in I \} = \{ \glssymbol{elementarereignis} \in \eraum  : \exists j \in I \text{ mit } \glssymbol{elementarereignis} \in \erga_j \}$
        \item[(b)] $\underset{i \in I}{\bigcap} \erga_i = \{ \glssymbol{elementarereignis} \in \eraum  : \glssymbol{elementarereignis} \in \erga_i \text{ für alle } i \in I\} = \{ \glssymbol{elementarereignis} \in \eraum  : \glssymbol{elementarereignis} \in \erga_i \forall j \in I \} $
    \end{itemize}
\end{defi}

\begin{lem}
    \textbf{Rechenregeln für Mengen und Ereignisse}  \txc{(2025, Skript 26 - Lemma 2.3)} \\
    Seien $\erga,\ergb,\ergc \subset \eraum $.
    \begin{figure}[H]
        \centering
        \begin{subfigure}[t]{0.45\textwidth}
            \centering
            \begin{tikzpicture}
                % Circle with label at 225 deg.
                \node[draw,
                    circle,
                    minimum size =3cm,
                    fill=accdarklight,
                ] (circle1) at (0,0){};

                % Circle with label at 45 deg.
                \node[draw,
                circle,
                minimum size =2cm,
                fill=accdarklight!80,
                label={280:$\ergc$}] (circle2) at (0,0.3){};

                \node[draw,
                circle,
                minimum size =1cm,
                fill=accdarklight!50,
                label={225:$\ergb$}] (circle3) at (0.4,0.4){};
                \node[black] at (circle3.center){$\erga$};
            \end{tikzpicture}
            \caption{Transitivität: $\erga\subset \ergb, \quad \ergb \subset \ergc \Rightarrow \erga \subset \ergc$}
            \label{fig:Transitivitaet}
        \end{subfigure}
        \hfill
        \begin{subfigure}[t]{0.45\textwidth}
            \centering
            \begin{tikzpicture}[thick,
                set/.style = {circle,
                minimum size = 3cm,
                fill=accdarklight!80}]
                % Set A
                \node[set,label={135:$\erga$}] (A) at (0,0) {};

                % Set B
                \node[set,label={45:$\ergb$}] (B) at (1.8,0) {};

                % Set C
                \node[set,label=$\ergc$] (C) at (0.9,1.5) {};

                % Intersection
                \begin{scope}
                    \clip (0,0) circle(1.5cm);
                    \clip (1.8,0) circle(1.5cm);
                    \clip (0.9,1.5) circle(1.5cm);
                    \fill[accdarklight!50](0,0) circle(1.5cm);
                \end{scope}

                % Circles outline
                \draw (0,0) circle(1.5cm);
                \draw (1.8,0) circle(1.5cm);
                \draw (0.9,1.5) circle(1.5cm);

            \end{tikzpicture}
            \caption{Distributivgesetz: $\erga \cap (\ergb \cup \ergc) = (\erga \cap \ergb) \cup (\erga \cap \ergc)$ und
                $ \erga \cup (\ergb \cap \ergc) = (\erga \cup \ergb) \cap (\erga \cup \ergc)$}
            \label{fig:Distributivgesetz}
        \end{subfigure}\label{fig:figure}
    \end{figure}
    \begin{itemize}
        \item Kommutativitätsgesetz: $\erga \cup \ergb = \ergb \cup \erga$ und $\erga \cap \ergb = \ergb \cap \erga$
        \item Assoziativgesetz:$(\erga \cup \ergb) \cup \ergc = \erga \cup (\ergb \cup \ergc) $ und $(\erga \cap \ergb) \cap \ergc = \erga \cap (\ergb \cap \ergc) $
        \item de Morgansche Regeln: $(\erga \cup \ergb)^c = \erga^c \cap \ergb^c$ und  $(\erga \cap \ergb)^c = \erga^c \cup \ergb^c$
    \end{itemize}
    \paragraph{Beweis.} Skript 26 Seite 20 % TODO
\end{lem}

\karkar{def_komplement}
\karkar{def_vereinigung_schnitt_abzaehlbar}
\karkar{lem_rechenregeln_mengen}

\subsection{Disjunktheit}
\begin{defi}
    $\erga, \ergb, \erga_1, \erga_2, \dots \subset \eraum $ \txc{(2025)}
    \begin{itemize}
        \item[(a)] $\erga$ und $\ergb$ heißen \underline{disjunkt} , wenn $\erga \cap \ergb = \emptyset$
        \begin{figure}[H]
            \centering
            \begin{tikzpicture}
                \node[draw,
                circle,
                minimum size =2cm,
                fill=accdarklight!50,
                label={225:$\erga$}] (circle1) at (0,0){};
                \node[draw,
                    circle,
                    minimum size =0.5cm,
                    fill=accdarklight!80,
                ] (circle1) at (0,0.6){};
                \node[draw,
                    circle,
                    minimum size =0.5cm,
                    fill=accdarklight!30,
                ] (circle1) at (0,-0.5){};
                \node[draw,
                    circle,
                    minimum size =0.5cm,
                    fill=accdarklight!70,
                ] (circle1) at (-0.3,0){};

                % Circle with label at 45 deg.
                \node[draw,
                circle,
                minimum size =2cm,
                fill=urgent!50,
                label={45:$\ergb$}] (circle2) at (3,0){};
                \node[draw,
                    circle,
                    minimum size =0.5cm,
                    fill=urgent!20,
                ] (circle2) at (3,-0.6){};
                \node[draw,
                    circle,
                    minimum size =0.5cm,
                    fill=urgent!70,
                ] (circle2) at (3,0.5){};
                \node[draw,
                    circle,
                    minimum size =0.5cm,
                    fill=urgent!90,
                ] (circle2) at (2.4,0){};
                \node[draw,
                    circle,
                    minimum size =0.5cm,
                    fill=urgent!20,
                ] (circle2) at (3.6,0){};
            \end{tikzpicture}
            \caption{Disunkte Mengen.}
            \label{fig:disjunkt}
        \end{figure}
        \item[] Sei $I$ eine beliebige Indexmenge. Für jedes $i \in I$ sei eine Menge $A_i \subseteq \eraum$ gegeben.
        \item[(b)] $\erga_1, \erga_2, \dots$  heißen \underline{disjunkt} , wenn $\underset{i}{\bigcap} \erga_i = \emptyset$ \txc{(Alle Mengen liegen getrennt voneinander)}
        \begin{figure}[H]
            \centering
            \begin{tikzpicture}
                % Set A
                \node [circle,
                fill=accdarklight!50,
                minimum size =3cm,
                label={135:$\erga_1$}] (A) at (0,0){};
                \node [circle,
                fill=accdarklight!50,
                minimum size =3cm,
                label={45:$\erga_3$}] (C) at (3.6,0){};

                \node [circle,
                fill=accdarklight!50,
                minimum size =3cm,
                label={90:$\erga_2$}] (B) at (1.8,0){};
                % Circles outline
                \draw[black,thick] (0,0) circle(1.5cm);
                \draw[black,thick] (1.8,0) circle(1.5cm);
                \draw[black,thick] (3.6,0) circle(1.5cm);
            \end{tikzpicture}
            \caption{Disjunkte Mengen.}
            \label{fig:DisjunkteMengen}
        \end{figure}
        \item[(c)] $\erga_1, \erga_2, \dots $ heißen paarweise disjunkt, wenn $\erga_i \cap \erga_j \neq \emptyset \quad \forall i,j $ mit $i \neq j$
        \begin{figure}[H]
            \centering
            \begin{tikzpicture}
                % Set A
                \node [circle,
                fill=accdarklight!50,
                minimum size =3cm,
                label={135:$\erga_1$}] (A) at (0,0){};
                \node [circle,
                fill=accdarklight!50,
                minimum size =3cm,
                label={45:$\erga_3$}] (C) at (7,0){};

                \node [circle,
                fill=accdarklight!50,
                minimum size =3cm,
                label={90:$\erga_2$}] (B) at (3.5,0){};
                % Circles outline
                \draw[black,thick] (0,0) circle(1.5cm);
                \draw[black,thick] (3.5,0) circle(1.5cm);
                \draw[black,thick] (7,0) circle(1.5cm);
            \end{tikzpicture}
            \caption{Disjunkte Mengenfamilie; paarweise disjunkte Mengen.}
            \label{fig:DisjunkteMengenfamilie}
        \end{figure}
        \item[(d)] Sind $\erga_1, \erga_2, \dots $ paarweis disjunkt, so schreiben wir $\underset{i}{\dot\cup} \erga_i$ für die Vereinigung  $\underset{i}{\cup} \erga_i$.
        \item Wir sprechen von einer \underline{disjunkten Vereinigung}
    \end{itemize}
\end{defi}

\begin{bem}
    \txc{(2025)}\\
    paarweise disjunkt $\Rightarrow$ disjunkt \\
    $\underset{i}{\cup} \erga_i \subset \erga_1 \cap \erga_2 = \emptyset$
\end{bem}

\begin{bsp} \textbf{Disjunkte und paarweise disjunkte Mengen beim Würfeln} \txc{(Skript 26 - Beispiel 2.5)} \\
% TODO das Beispiel aus 2025 fehlt
    Sei $\eraum = \{1,...,6\}$
    \begin{itemize}
        \item Setze $A = \{2, 4, 6\}, B = \{3, 6\}$ und $C = \{1\}$. Dann ist $A \cap B = \{6\}̸ = \emptyset$, also sind A
        und B nicht disjunkt, während $ A \cap C = \emptyset$ impliziert, daß A und C disjunkt sind.
        \item Seien $A_1 = \{1,3\}$, $A_2 = \{3, 5\}$ und $A_3 = \{1, 5\}$. Dann sind $A_1, A_2$ und $A_3$ disjunkt,
        aber nicht paarweise disjunkt, weil $A_1 \cap A_2 \cap A_3 = \emptyset$, aber $A_1 \cap A_2 = \{3\}̸ = \emptyset$.
    \end{itemize}
\end{bsp}

\begin{lem}
    \textbf{Relation zwischen den Konzepten ``disjunkt'' und ``paarweise disjunkt''}
    \txc{(Skript 26 - Lemma 2.6)} \\
    Sei I eine beliebige Indexmenge.
    Für jedes $i \in I$ sei eine Menge $A_i \subseteq \eraum$ gegeben.
    \begin{enumerate}
        \item paarweise disjunkt $\Rightarrow$ disjunkt: \\ $A_i \cap A_j = \emptyset \quad \forall i,j \in I$ mit $i \neq j \Rightarrow \bigcap_{i \in I} A_i = \emptyset$
        \item disjunkt $\nRightarrow$ paarweise disjunkt. \\Es gibt Beispiele von disjunkten Mengen, die nicht
        paarweise disjunkt sind.
    \end{enumerate}
    \paragraph{Beweis.} Skript 26 Seite 22 % TODO 
\end{lem}

\begin{bsp}
    \textbf{Disjunkte und paarweise disjunkte Mengen} \txo{(Skript 26 - Beispiel 2.7)}
    Wir wählen
    \begin{align*}
        A_1 & =\{1,2,3,4,\dots\}, \\
        A_2 & =\{2,3,4,\dots\}, \\
        A_3 & =\{3,4,\dots\}, \\
        \vdots \\
        A_k & =\{k,k+1,\dots\}=\{j\in\mathbb{N}:j\ge k\}.
    \end{align*}
    Die Mengen \(A_i\), \(i\in\mathbb{N}\), sind \emph{nicht paarweise disjunkt}:\\
    Für jede Wahl von \(i,j\in\mathbb{N}\) mit \( i < j \) gilt
    \[
        A_i\cap A_j=A_j\neq\emptyset.
    \]
    Dennoch sind die Mengen \(A_i\), \(i\in\mathbb{N}\) \textit{disjunkt}
    \[
        \bigcap_{i\in\mathbb{N}}A_i=\emptyset.
    \]
    Zum Beweis nehmen wir an, es gäbe ein \(\ell\in\mathbb{N}\) mit
    \[
        \ell\in\bigcap_{i\in\mathbb{N}}A_i.
    \]
    Dann müsste \(\ell\in A_i\) für alle \(i\in\mathbb{N}\) gelten. Insbesondere müsste
    \[
        \ell\in A_{\ell+1}=\{\ell+1,\ell+2,\dots\}
    \]
    sein. Dies ist jedoch unmöglich. Also existiert kein solches
    \(\ell\), und damit folgt
    \[
        \bigcap_{i\in\mathbb{N}}A_i=\emptyset.
    \]
\end{bsp}

\karkar{def_disjunkt_zwei_mengen}
\karkar{def_disjunkt_mengenfamilie}
\karkar{def_paarweise_disjunkt}
\karkar{def_disjunkte_vereinigung_notation}
\karkar{bem_paarweise_disjunkt_impliziert_disjunkt}
\karkar{bsp_disjunkt_wuerfeln}
\karkar{lem_disjunkt_nicht_paarweise_disjunkt}
\karkar{bsp_disjunkt_nicht_paarweise_Ak}


\section{Erste Rechenregeln für Wahrscheinlichkeiten}\label{sec:erste-rechenregeln-fur-wahrscheinlichkeiten}
\begin{satz}
    \txc{(2025)} \\
    Seien $\erga, \ergb, \erga_1, \erga_2, \dots \subset \eraum $ (und Intervalle, falls $\pe$ durch \dichte gegeben ist)
    \begin{itemize}
        \item[(a)] Additivität: Seien $\erga_1, \erga_2, \dots$ paarweise disjunkt, dann gilt:
        \[\pe(\underset{i}{\dot\cup} \erga_i) = \sum_i \pe(\erga_i)\]
        Gilt auch für \Glspl{dichte}! \\
        \textbf{Beweis.}  \\
        \begin{align*}
            \pe(\underset{i}{\dot\cup} \erga_i)&  = \pe(\{\glssymbol{elementarereignis} \in \eraum : \exists i \text{ mit } \glssymbol{elementarereignis} \in \erga_i\}) \\
            & = \sum_i \txo{\underbrace{\textcolor{black}{\sum_{\glssymbol{elementarereignis} \in \erga_i}}}_{\pe(\erga_i)}} \glssymbol{wfk} (\glssymbol{elementarereignis}) \\
            & = \sum_i \pe(\erga_i)
        \end{align*}
        \orangenote{Wir haben die Summe/Reihe umsortiert}

        \item[(b)] Komplementärwahrscheinlichkeit (``Gegenwahrscheinlichkeit'')
        \[\pe(\erga^c) = \pe(\eraum  \backslash \erga) = 1 - \pe(\erga)\]
        \textbf{Beweis.}
        \begin{align*}
            \eraum  =  \erga \dot\cup \erga^c & \overset{(a)}{\Rightarrow} 1 \overset{\txo{Normierung}}{=} \pe(\eraum ) = \pe( \erga \dot\cup \erga^c) \overset{(a)}{=}\pe(\erga) + \pe(\erga^c) \\
            & \Rightarrow \pe(\erga^c) = 1-  \pe(\erga)
        \end{align*}
        \item[(c)] Monotonie: \[\erga \subset \ergb \Rightarrow \pe(\erga) \leq \pe(\ergb)\]
        \textbf{Beweis}
        \begin{align*}
            \ergb & = \erga \dot\cup (\ergb\backslash \erga) \\
            \pe(\ergb) & = \pe(\erga) + \txo{
                \underbrace{\textcolor{black}{\pe(\ergb \backslash \erga) }}_{\geq 0}
            } \geq \pe(\erga)
        \end{align*}
        \item[(d)] $\erga \subset \ergb\Rightarrow \pe(\erga) = \pe(\ergb) - \pe(\ergb \backslash \erga) $
        \item[(e)] Wahrscheinlichkeit von Vereinigungen: \[\pe(\erga \cup \ergb) = \pe(\erga) + \pe(\ergb) - \pe(\erga \cap \ergb)\] \\
        \textbf{Beweis.}
        \begin{align*}
            \erga \cup \ergb & = \erga \dot\cup [\ergb \backslash (\erga \cap \ergb)] \\
            \pe(\erga \cup \ergb) & = \pe(\erga) +
            \txo{
                \underbrace{
                    \textcolor{black}{\pe(\ergb \backslash (\erga \cap \ergb) )}
                }_{ \overset{(d)}{=}\pe(\ergb) - \pe(\erga \cap \ergb) }
            } \\
            & = \pe(\erga) + \pe(\ergb) - \pe(\erga \cap \ergb)
        \end{align*}
    \end{itemize}
    Weitere Beweise Skript 2026 ab Seite 23. % TODO

\end{satz}
\begin{bem}
    \txc{(2025)} \\
    Sind $\erga, \ergb$ disjunkt, so gilt nach (a): $\pe(\erga \dot\cup \ergb)  = \pe(\erga) + \pe(\ergb) $. \\
    Sind $\erga, \ergb$ nicht disjunkt, so gilt nach (e): $\pe(\erga \cup \ergb) = \pe(\erga) + \pe(\ergb) - \pe(\erga \cap \ergb)$. \\
    Sind $\erga, \ergb, \ergc$ paarweise disjunkt, so gilt nach (a): $\pe(\erga \cup \ergb \cup \ergc) \overset{(a)}{=} \pe(\erga) + \pe(\ergb) + \pe(\ergc) $. \\ \\
    Was passiert, wenn $\erga,\ergb,\ergc$ nicht paarweise disjunkt sind? \begin{align*}
                                                                              \pe(\erga \cup \ergb \cup \ergc) & = \pe(\erga \cup (\ergb \cup \ergc)) \\
                                                                              & = \pe(\erga) \pe( \cup \ergb \cup \ergc) -  \pe(\erga \cap (\ergb \cup \ergc)) \\
                                                                              & = \pe(\erga) + [\pe(\ergb) + \pe(\ergc) - \pe(\ergb \cap \ergc)] - \pe((\erga \cap \ergb) \cup (\erga \cap \ergc))\\
                                                                              & = \pe(\erga) + \pe(\ergb) + \pe(\ergc)  & \text{Mengen }  \\
                                                                              & \quad - \pe(\erga \cap \ergb) - \pe(\erga \cap \ergc) - \pe(\ergb \cap \ergc) & \text{Paare von Mengen }  \\
                                                                              & \quad + \pe(
                                                                              \txo{
                                                                                  \underbrace{
                                                                                      \textcolor{black}{(\erga \cap \ergb) \cap (\erga \cap \ergc)}
                                                                                  }_{= \erga \cap \ergb \cap \ergc}
                                                                              }
                                                                              ) & \text{Tripel von Mengen }
    \end{align*}
\end{bem}

\karkar{satz_additivitaet}
\karkar{satz_komplementaerwahrscheinlichkeit}
\karkar{satz_monotonie}
\karkar{satz_differenz_teilmenge}
\karkar{satz_vereinigung_zwei_mengen}
\karkar{bem_zusammenfassung_vereinigungsregeln}

\subsection{Das Ein- und Ausschlußprinzip}
\begin{satz}
    \textbf{Ein- und Ausschlußprinzip für drei Ereignisse} \txc{(Skript 26 - Lemma 2.9)}
    \begin{align*}
        P(A \cup B \cup C)
        &= P(A) + P(B) + P(C) && \text{(Mengen)}\\
        &\quad - P(A \cap B) - P(A \cap C) - P(B \cap C)
        && \text{(Paare von Mengen)}\\
        &\quad + P(A \cap B \cap C)
        && \text{(Tripel von Mengen)}
    \end{align*}
    \paragraph{Beweis.}
    \begin{align*}
        P(A \cup B \cup C)
        &= P((A \cup B) \cup C)\\
        &= P(A \cup B) + P(C) - P((A \cup B) \cap C)\\
        &= P(A \cup B) + P(C) - P((A \cap C) \cup (B \cap C))\\
        &= P(A) + P(B) - P(A \cap B) + P(C) \\
        &\quad - \Bigl[
            P(A \cap C) + P(B \cap C) - P((A \cap C) \cap (B \cap C))
            \Bigr]\\
        &= P(A) + P(B) + P(C) - P(A \cap B) - P(A \cap C) \\
        &\quad - P(B \cap C) + P(A \cap B \cap C).
    \end{align*}
\end{satz}

\begin{satz}
    \textbf{Ein- und Ausschlußprinzip für n Ereignisse}  \txc{(2025, Skript 26 - Theorem 2.10)} \\
    \[\erga_2,\dots, \erga_n \subset \eraum  \]
    \[\pe\left( \bigcup^n_{i=1} \erga_i \right) = \sum^n_{k = 1} (-1)^{k+1} \sum_{a \leq i_1 \leq \dots \leq i_k \leq k} \pe(\erga_{i_1} \cap \dots \cap \erga_{i_k}) \]
    \paragraph{Beweis.} Die Fälle n = 2 und n = 3 haben wir bereits behandelt. Der Beweis für beliebiges n kann mit Hilfe
    vollständiger Induktion geführt werden.
\end{satz}

\begin{bsp}
    \textbf{Werfen einer Sechs bei n-maligem Würfeln}  \txc{(Skript 26 - Beispiel 2.11)} \\
    Wie groß ist die Wahrscheinlichkeit, bei n-maligem Würfeln mit einem fairen Würfel mindestens eine
    Sechs zu werfen? \\
    Siehe Skript 2026 Seite 25
\end{bsp}

% TODO Beispiel 2.12 (Knobelaufgabe)

\karkar{satz_einschluss_ausschluss_drei}
\karkar{satz_einschluss_ausschluss_n}

\newpage

